% 13
\begin{frame}[t]{与医院合作的流程与关卡}
\vspace{.1cm}
\centering
\begin{tikzpicture}[sg/.style={box,align=left,inner sep=7pt,text width=2.2cm,minimum height=2.6cm,font=\scriptsize},
  gate/.style={diamond,draw=DarkBlue,fill=white,line width=.8pt,inner sep=2.4pt},
  gl/.style={font=\tiny,text=DarkBlue,align=center,text width=.85cm,anchor=north,inner sep=1pt}]
 \node[sg,fill=PaleAmber] (s1) at (0,0) {{\small\bfseries\color{DarkBlue}\faHandshake\ 合作前}\\[.12cm]{\color{Muted}$\bullet$}\ 伦理批件是否覆盖 AI 分析与二次利用\\[.04cm]{\color{Muted}$\bullet$}\ 人类遗传资源相关审批\\[.04cm]{\color{Muted}$\bullet$}\ 数据使用协议与成果归属};
 \node[sg] (s2) at (3.6,0) {{\small\bfseries\color{DarkBlue}\faIcon{shield-alt}\ 数据准备}\\[.12cm]{\color{Muted}$\bullet$}\ 脱敏标准与责任人\\[.04cm]{\color{Muted}$\bullet$}\ 数据能否出院区\\[.04cm]{\color{Muted}$\bullet$}\ 院内 GPU 或导出流程};
 \node[sg] (s3) at (7.2,0) {{\small\bfseries\color{DarkBlue}\faClipboardCheck\ 实验与核验}\\[.12cm]{\color{Muted}$\bullet$}\ 谁提供标签或判读\\[.04cm]{\color{Muted}$\bullet$}\ 结果如何回到临床验证\\[.04cm]{\color{Muted}$\bullet$}\ 合作者的时间投入};
 \node[sg] (s4) at (10.8,0) {{\small\bfseries\color{DarkBlue}\faFlagCheckered\ 成果}\\[.12cm]{\color{Muted}$\bullet$}\ 发表前审查流程\\[.04cm]{\color{Muted}$\bullet$}\ 署名与知识产权\\[.04cm]{\color{Muted}$\bullet$}\ 数据留存与销毁};
 \draw[flow] (s1)--(s2);\draw[flow] (s2)--(s3);\draw[flow] (s3)--(s4);
 \node[gate] (g1) at (1.8,0) {};\node[gate] (g2) at (5.4,0) {};\node[gate] (g3) at (9.0,0) {};
 \node[gl] at ([yshift=-.14cm]g1.south) {批件\\到位};
 \node[gl] at ([yshift=-.14cm]g2.south) {数据流向\\确定};
 \node[gl] at ([yshift=-.14cm]g3.south) {核验\\通过};
 \node[font=\scriptsize,text=Muted,anchor=north,align=center,text width=3cm] at ([yshift=-.15cm]s1.south) {\faHourglassHalf\enspace 常需数月\\与起步阶段并行启动};
\end{tikzpicture}
\vspace{.2cm}
\remark{\ico{\faExclamationTriangle}\enspace 菱形为关卡；未通过前不把医院数据纳入任何实验，期间用公开数据推进起步阶段。}
\source{合作前清单，不构成法律意见；具体要求以医院伦理委员会与主管部门的规定为准。}
\end{frame}
